\section{Análise de Erros e a Natureza Morfológica das Classes}

A análise qualitativa das classificações incorretas revelou que o sistema não falha de forma aleatória, mas sim em conformidade com as fronteiras morfológicas da fibra de algodão. As confusões concentram-se predominantemente entre classes adjacentes (como Tipo 3 e Tipo 4), o que reflete a subjetividade intrínseca enfrentada por classificadores humanos.

O fato de a classe \textbf{41.4} ser frequentemente confundida com a \textbf{31.2} evidencia que existem limites físicos de discriminabilidade visual que até mesmo modelos de Deep Learning encontram dificuldade em superar. Esta observação sugere que a classificação automatizada pode ser utilizada não apenas como um substituto, mas como uma ferramenta de apoio à decisão, fornecendo um grau de incerteza (probabilidade) que permite ao classificador humano intervir apenas nos casos de fronteira, reduzindo drasticamente o tempo total de processamento no emblocamento industrial.
